\chapter{Related Work}

This chapter surveys the research landscape that underpins the proposed framework for target-driven, memory-attention-enhanced humanoid parkour. We organize the review into four thematic areas: humanoid locomotion and parkour (Section~2.1), memory-augmented policies (Section~2.2), attention and mixture-of-experts mechanisms in robot learning (Section~2.3), and adversarial motion priors (Section~2.4). Section~2.5 summarizes the key limitations of existing work and identifies the research gap that this thesis aims to fill.

%% ============================================================
\section{Humanoid Locomotion and Parkour}
\label{sec:rw_locomotion}
%% ============================================================

\subsection{Background and Motivation}

Humanoid robots, by virtue of their anthropomorphic morphology, hold the promise of operating in environments originally designed for humans---climbing stairs, traversing uneven terrain, squeezing through narrow passages, and performing dynamic maneuvers that require whole-body coordination. The quest to endow humanoid robots with agile locomotion capabilities has thus become one of the central challenges in robotics research~\cite{gu2024humanoidlocomotion}. Unlike wheeled or tracked platforms, legged systems must continuously manage balance, contact sequencing, and terrain interaction, all while respecting the underactuated nature of floating-base dynamics. These challenges are amplified in parkour scenarios, where the robot must execute highly dynamic, acrobatic skills---such as vaulting, climbing, leaping, and crawling---in rapid succession across complex, often unstructured environments.

The term ``parkour'' in the robotics context refers to the ability of a legged robot to navigate through obstacle-laden environments using a repertoire of dynamic whole-body skills, drawing inspiration from the human discipline of traversing urban and natural obstacles with speed and efficiency. While traditional locomotion research has focused on steady-state walking and running on flat or mildly uneven terrain, parkour demands far more aggressive dynamics, diverse contact modes (hands, knees, torso), and perceptive awareness of the surrounding geometry. As surveyed by Gu et al.~\cite{gu2024humanoidlocomotion}, the field of humanoid locomotion has evolved from model-based control approaches toward learning-based methods that can handle the complexity of whole-body coordination, and parkour represents the frontier of this evolution.

\subsection{Reinforcement Learning for Legged Locomotion}

Reinforcement learning (RL) has emerged as the dominant paradigm for training locomotion controllers, owing to its ability to discover complex control policies through trial-and-error interaction with simulated environments. The foundational algorithm Proximal Policy Optimization (PPO)~\cite{schulman2017ppo} has become the de facto standard for continuous-control RL in robotics due to its stability, scalability, and ease of implementation. PPO's clipped surrogate objective ensures monotonic policy improvement, making it well-suited for the high-dimensional action spaces of humanoid robots.

Rudin et al.~\cite{rudin2016walk} demonstrated that massively parallel simulation enables training quadrupedal locomotion policies in mere minutes, a breakthrough that dramatically reduced the computational barrier to RL-based locomotion. By leveraging thousands of simulated environments running concurrently, their approach showed that sample-efficient training is possible even for high-dimensional legged systems. This insight has been foundational for subsequent work on humanoid locomotion, where the action space is even larger and the dynamics more complex.

Xie et al.~\cite{xie2024leggedlocomotion} provided a comprehensive survey of reinforcement learning for legged robot locomotion, identifying key design choices including reward shaping, observation space design, domain randomization strategies, and sim-to-real transfer techniques. They highlighted that the success of RL-based locomotion hinges on careful engineering of the training pipeline---particularly the formulation of reward functions that balance tracking accuracy, energy efficiency, and robustness to perturbations. Their analysis underscored that while RL has achieved impressive results on quadrupedal platforms, extending these methods to humanoid robots with more degrees of freedom and more complex contact dynamics remains an open challenge.

Margolis et al.~\cite{margolis2024rapidlocomotion} proposed a rapid locomotion framework that combines reinforcement learning with a set of carefully designed reward components to achieve agile running and turning behaviors on quadrupedal robots. Their approach demonstrated that simple reward structures, when combined with appropriate domain randomization, can produce policies that transfer effectively to real hardware. The principles established in this work---modular reward design, extensive randomization, and efficient simulation---have been widely adopted in subsequent humanoid locomotion research.

DreamWaQ~\cite{nahraendra2024dreamwaq} introduced the concept of implicit terrain imagination for robust quadrupedal locomotion. By training a latent representation that implicitly encodes terrain properties, DreamWaQ enables the policy to make informed locomotion decisions even when explicit terrain observations are noisy or incomplete. The key insight is that the policy's internal representation can serve as a ``mental model'' of the terrain, allowing it to adapt its gait and foothold selection without requiring explicit terrain reconstruction. This idea of implicit terrain understanding has influenced subsequent work on perceptive humanoid locomotion, where the complexity of the terrain demands richer internal representations.

Ai et al.~\cite{ai2024embodimentscaling} investigated embodiment scaling laws in robot locomotion, studying how policy performance scales with model capacity, training data, and embodiment complexity. Their findings suggested that larger models and more diverse training environments lead to better generalization, but that the relationship is not purely monotonic---certain architectural choices and training strategies can yield disproportionate improvements. This work provides important context for understanding the trade-offs involved in designing humanoid locomotion policies with varying levels of model complexity.

\subsection{Perceptive Locomotion on Challenging Terrains}

A critical capability for real-world deployment is the ability to perceive and adapt to terrain geometry. Blind locomotion policies, while effective on flat ground, fail catastrophically when confronted with stairs, gaps, slopes, or sparse footholds. The integration of perceptual information---typically from depth cameras or LiDAR---into the locomotion control loop has thus become a major research direction.

Bank et al.~\cite{bank2024hybrid} proposed a hybrid autoencoder for robust heightmap generation from fused LiDAR and depth data, specifically designed for humanoid robot locomotion. Their method addresses the challenge of combining heterogeneous sensor modalities---LiDAR providing sparse but accurate long-range measurements and depth cameras providing dense but noisy short-range data---into a unified heightmap representation suitable for policy consumption. The hybrid autoencoder architecture learns to fuse these modalities in a way that is robust to sensor noise, occlusion, and varying point densities, producing heightmaps that enable more reliable foot placement decisions. This work highlights the importance of perceptual preprocessing in the locomotion pipeline, as the quality of the terrain representation directly impacts the policy's ability to make safe and efficient locomotion decisions.

RPL~\cite{zhang2024rpl} presented a framework for learning robust humanoid perceptive locomotion on challenging terrains. Their approach combines a privileged information encoder that processes terrain heightmaps during training with a student policy that relies on proprioceptive history and exteroceptive observations during deployment. The key contribution is a carefully designed curriculum that progressively increases terrain difficulty, enabling the policy to develop robust locomotion strategies that generalize to unseen terrain configurations. RPL demonstrated successful sim-to-real transfer on a Unitree G1 humanoid robot navigating slopes, stairs, and rough terrain, establishing a strong baseline for perceptive humanoid locomotion.

Miki et al.~\cite{miki2024confinedspaces} addressed the problem of learning to walk in confined spaces using 3D representations. Unlike open-terrain locomotion where the primary concern is ground geometry, confined-space navigation requires awareness of overhead and lateral obstacles---walls, ceilings, and narrow passages that constrain the robot's whole-body posture. Their method employs a 3D voxel-based representation of the surrounding environment, enabling the policy to plan body-level maneuvers such as crouching, side-stepping, and crawling. This work is particularly relevant to parkour scenarios, where the robot must frequently navigate through tight spaces such as tunnels, gaps, and under obstacles.

Zhang et al.~\cite{zhang2024riskyterrain} focused on learning agile locomotion on risky terrains, where a single misstep can lead to catastrophic failure. Their approach emphasizes the importance of risk-aware foot placement, training the policy to prioritize safe footholds while maintaining forward progress. By incorporating a risk metric into the reward function and using a terrain difficulty curriculum, their method achieves reliable locomotion on terrains with sparse footholds, narrow ridges, and large gaps. The concept of risk-aware locomotion is central to parkour, where the robot must balance aggressive maneuver execution with safety constraints.

\subsection{Humanoid and Legged Parkour}

The extension from steady-state locomotion to dynamic parkour represents a qualitative leap in capability. Parkour requires the robot to chain multiple dynamic skills---jumping, climbing, vaulting, crawling---in rapid succession, adapting to the specific geometry of each obstacle. This subsection reviews the major works in legged and humanoid parkour.

Cheng et al.~\cite{cheng2024extremeparkour} introduced Extreme Parkour for legged robots, demonstrating that quadrupedal robots can execute highly dynamic maneuvers including high jumps, long jumps, and tilting ramps. Their approach trains separate skills for each maneuver type and uses a high-level policy to select and parameterize skills based on the perceived terrain. The key insight is that decomposing parkour into a library of specialized skills, each trained with a tailored reward structure, is more effective than attempting to learn a single monolithic parkour policy. This skill-chaining paradigm has been widely adopted in subsequent humanoid parkour work.

Robot Parkour Learning~\cite{zhuang2024robotparkour} extended the parkour framework to quadrupedal robots with a focus on learning a diverse repertoire of parkour skills through reinforcement learning. Their method employs a two-level hierarchy: a low-level controller that executes individual skills and a high-level planner that sequences skills based on terrain perception. The training procedure uses extensive domain randomization and a progressive difficulty curriculum to ensure that the learned skills are robust to sim-to-real transfer gaps. This work demonstrated that quadrupedal robots can perform parkour maneuvers such as leaping over gaps, climbing onto platforms, and crawling under barriers with reliable real-world performance.

Humanoid Parkour Learning~\cite{zhuang2024humanoidparkour} brought the parkour paradigm to humanoid robots, addressing the additional challenges posed by the humanoid morphology---more degrees of freedom, a higher center of mass, and the need for whole-body coordination including upper-body movements. Their approach trains a unified policy that can execute diverse parkour skills including jumping, climbing, and crawling, using a single reinforcement learning framework. A key contribution is the design of a reward function that encourages dynamic whole-body motion while maintaining balance and respecting joint limits. The method was validated on a Unitree H1 humanoid robot, demonstrating successful execution of parkour courses in both simulation and the real world.

Deep Whole-body Parkour~\cite{zhuang2024deepwholebody} further advanced humanoid parkour by proposing an end-to-end learning framework that directly maps terrain observations to whole-body actions, eliminating the need for explicit skill decomposition. Their method trains a single deep policy that learns to perform a wide variety of parkour maneuvers---from precision jumping to wall-climbing to rolling---by leveraging a large-scale simulated training environment with diverse obstacle configurations. The policy architecture incorporates a terrain encoder that processes heightmap observations and a proprioceptive encoder that processes the robot's state history, with both streams fused before the action head. This work demonstrated that end-to-end training can produce more fluid and adaptive parkour behaviors compared to skill-chaining approaches, at the cost of requiring significantly more training data and compute.

Wu et al.~\cite{wu2024perceptiveparkour} proposed Perceptive Humanoid Parkour, which chains dynamic human skills via motion matching. Their approach combines motion capture data of human parkour athletes with reinforcement learning to train a humanoid robot to replicate dynamic parkour behaviors. The motion matching component retrieves the most appropriate reference motion from a library based on the current terrain and robot state, while the RL policy tracks this reference while adapting to the specific dynamics of the robot and terrain. This hybrid approach leverages the richness of human motion data while maintaining the adaptability of learned policies, achieving parkour behaviors that are both natural-looking and physically feasible.

Hiking in the Wild~\cite{zhu2024hiking} presented a scalable perceptive parkour framework for humanoids that emphasizes generalization to diverse, unstructured outdoor environments. Unlike previous works that primarily evaluate on structured obstacle courses, this framework is designed to handle the variability of natural terrain---loose rocks, steep slopes, dense vegetation, and irregular surfaces. The method employs a perceptive encoder that processes depth observations into a compact terrain representation, combined with a locomotion policy trained on a large-scale terrain curriculum. A key contribution is the use of automated terrain generation that produces diverse training environments, enabling the policy to develop robust locomotion strategies that transfer to real outdoor settings.

TTT-Parkour~\cite{zhu2024tttparkour} introduced rapid test-time training for perceptive robot parkour, addressing the challenge of generalizing to terrain types not seen during training. Their method leverages the test-time training paradigm: at deployment, the policy rapidly adapts to the local terrain by performing a few gradient steps on a self-supervised objective. This enables the policy to quickly calibrate its terrain understanding to novel environments without requiring additional supervised data or human intervention. TTT-Parkour demonstrated that test-time adaptation can significantly improve parkour performance on out-of-distribution terrains, suggesting that static policies may be insufficient for truly general parkour and that adaptive mechanisms are necessary.

PIE~\cite{luo2024pie} proposed a parkour framework with implicit-explicit learning for legged robots. The implicit component learns a latent terrain representation that captures the essential features of the environment, while the explicit component directly processes terrain heightmaps. By combining both streams, PIE achieves robust perceptive locomotion that is both data-efficient and generalizable. The implicit representation enables the policy to make informed decisions even when explicit terrain observations are degraded, while the explicit representation provides precise geometric information for foot placement and obstacle avoidance. This dual-stream approach has influenced subsequent work on multi-modal terrain processing for parkour.

\subsection{Specialized Parkour Skills and Platforms}

Beyond the general parkour frameworks discussed above, several works have focused on developing specialized skills that are essential components of a complete parkour repertoire.

APEX~\cite{wang2024apex} addressed the problem of learning adaptive high-platform traversal for humanoid robots. Climbing onto high platforms is one of the most challenging parkour skills, requiring precise coordination of arm pull-up, leg push-off, and body rotation. APEX trains a hierarchical policy where a high-level planner determines the approach trajectory and a low-level controller executes the climbing motion. The training employs a curriculum that starts with low platforms and progressively increases the height, enabling the policy to develop the strength and coordination needed for high-platform traversal. APEX demonstrated successful climbing onto platforms up to 1.5 times the robot's leg length on a Unitree H1 humanoid.

BeamDojo~\cite{wang2024beamdojo} focused on learning agile humanoid locomotion on sparse footholds, a critical capability for parkour scenarios where the robot must traverse narrow beams, stepping stones, or thin rails. The challenge lies in the precision required for foot placement---even small errors can lead to falls. BeamDojo employs a dual-stream observation architecture that processes both local foothold geometry and global terrain context, enabling the policy to plan precise stepping sequences. The training uses a specialized reward that heavily penalizes foot placement errors and a curriculum that gradually reduces foothold size and increases spacing. This work demonstrated reliable locomotion on beams as narrow as 10\,cm and stepping stones with gaps up to 40\,cm.

Walk the PLANC~\cite{dai2024walkplanc} proposed physics-guided reinforcement learning for agile humanoid locomotion on constrained footholds. Their method integrates physics-based priors---specifically, a model-predictive control (MPC) formulation that provides approximate solutions for foot placement and center-of-mass trajectory---into the RL training loop. The physics guidance serves as an inductive bias that accelerates learning and improves sample efficiency, while the RL component refines the policy beyond the accuracy limits of the simplified model. This hybrid approach achieved superior performance on narrow ridges, stepping stones, and discontinuous surfaces compared to pure RL methods.

Learning Visual Parkour from Generated Images~\cite{yu2024visualparkour} addressed the data scarcity problem in visual parkour by training policies using synthetically generated terrain images. Their method uses a generative model to produce diverse terrain images with associated depth maps, enabling large-scale training of visual parkour policies without requiring real-world data collection. The generated images span a wide distribution of terrain types, obstacles, and lighting conditions, providing the policy with exposure to scenarios that would be difficult or dangerous to encounter during physical training. This data augmentation strategy significantly improved the policy's ability to generalize to novel terrain configurations at deployment.

Let Humanoids Hike!~\cite{lin2024lethumanoids} proposed an integrative skill development framework for humanoid robots on complex hiking trails. Their approach decomposes the hiking task into a set of fundamental locomotion skills---walking, climbing, descending, and balancing---and trains each skill independently before integrating them through a high-level skill selector. The key insight is that complex trail navigation requires not only individual skill competence but also smooth transitions between skills, which is achieved through a transition-aware training procedure that explicitly optimizes for skill-switching stability. This work demonstrated that integrative skill development can produce more robust and natural hiking behaviors compared to monolithic policy approaches.

ASAP~\cite{he2024asap} focused on aligning simulation and real-world physics for learning agile humanoid whole-body skills. The sim-to-real gap---the discrepancy between simulated and real-world dynamics---is a fundamental challenge for RL-based locomotion, as policies trained in simulation often fail when deployed on physical hardware. ASAP proposes a systematic approach to closing this gap through a combination of system identification, dynamics randomization, and adaptive control. By carefully calibrating the simulation environment to match the real robot's dynamics and training the policy to be robust to residual discrepancies, ASAP achieved reliable transfer of agile whole-body skills including jumping, spinning, and waving from simulation to a real Unitree G1 humanoid.

FALCON~\cite{zhang2024falcon} extended the locomotion paradigm to loco-manipulation, learning force-adaptive humanoid control that coordinates locomotion with object interaction. While traditional parkour focuses on terrain traversal without object manipulation, many real-world scenarios require the robot to interact with objects---pushing doors, carrying tools, or manipulating levers---while navigating. FALCON trains a unified policy that simultaneously controls locomotion and manipulation, with a force-adaptive module that adjusts the robot's interaction force based on the object's physical properties. This work demonstrated that integrating locomotion and manipulation in a single policy can produce more efficient and robust behaviors compared to separate locomotion and manipulation controllers.

\subsection{Discussion and Limitations of Current Approaches}

The works reviewed in this section collectively demonstrate significant progress in humanoid locomotion and parkour, progressing from basic walking on flat terrain to dynamic whole-body parkour on complex obstacles. However, several important limitations remain.

First, most existing parkour policies are \emph{reactive} rather than \emph{predictive}: they respond to the current terrain observation without maintaining an internal model of the environment that supports long-horizon planning. This limits their ability to anticipate upcoming obstacles, prepare appropriate skills in advance, and chain skills smoothly over extended parkour courses. The implicit terrain imagination introduced by DreamWaQ~\cite{nahraendra2024dreamwaq} and the implicit-explicit learning of PIE~\cite{luo2024pie} represent steps toward predictive terrain understanding, but these approaches are primarily designed for quadrupedal locomotion and have not been extended to the full complexity of humanoid parkour with diverse skill repertoires.

Second, the \emph{target-driven} aspect of parkour---where the robot must navigate toward a specific goal while selecting appropriate skills---has received limited attention. Most existing frameworks assume that the robot follows a predefined course or receives direct velocity commands, rather than autonomously planning a path through the environment toward a distant target. The absence of target-driven planning limits the applicability of current parkour systems in scenarios where the robot must make high-level navigation decisions, such as choosing between alternative paths or determining the optimal sequence of skills to reach a goal.

Third, the \emph{memory} requirements of parkour have been largely overlooked. Parkour courses often span tens of meters, and the robot must remember previously observed terrain features---such as the location of a gap it passed earlier or the height of a wall it saw from a distance---to make informed decisions later in the course. Current policies typically rely on a short history of proprioceptive and exteroceptive observations, which may be insufficient for maintaining the spatial awareness needed for complex parkour navigation.

These limitations motivate the exploration of memory-augmented and attention-based policy architectures, which we review in the following sections.

%% ============================================================
\section{Memory-Augmented Policies}
\label{sec:rw_memory}
%% ============================================================

\subsection{The Need for Memory in Robot Control}

Standard reinforcement learning policies for locomotion are typically Markovian: they map the current observation (or a short history of recent observations) directly to an action. This Markovian assumption is reasonable for simple locomotion tasks such as walking on flat terrain, where the current observation contains sufficient information for decision-making. However, for complex tasks such as parkour, the Markovian assumption breaks down for several reasons.

First, the robot's sensors provide only a local, partial view of the environment. A depth camera mounted on the robot's head captures the terrain within a limited field of view and range, but terrain features outside this view---such as obstacles behind the robot or around corners---are invisible. To navigate effectively, the robot must accumulate observations over time to build a more complete understanding of the environment. This requires memory that extends beyond the immediate sensory horizon.

Second, many locomotion decisions depend on the history of the robot's interaction with the environment. For example, the appropriate gait for descending a slope depends on whether the robot just climbed a wall or walked across a flat surface, because the robot's momentum, joint configurations, and fatigue state differ in these scenarios. Similarly, the decision to jump across a gap may depend on whether the robot has already observed a safe landing surface on the other side during a previous pass. These history-dependent decisions require the policy to maintain and access information from previous time steps.

Third, in target-driven navigation, the robot must maintain a representation of its goal and its progress toward that goal. This requires persistent memory that stores the target location and updates the robot's estimated position relative to the target as it moves through the environment. Without such memory, the robot cannot make coherent navigation decisions over extended horizons.

The importance of memory in robot control has been recognized in several recent works, which we review below. These works can be broadly categorized into two approaches: (1) recurrent memory, where the policy maintains a hidden state that is updated at each time step through a recurrent neural network (RNN), and (2) external memory, where the policy reads from and writes to an explicit memory buffer that stores information over longer time scales.

\subsection{Recurrent Memory in Locomotion Policies}

The most common approach to incorporating memory into locomotion policies is through recurrent architectures, particularly Long Short-Term Memory (LSTM) networks and Gated Recurrent Units (GRUs). These networks maintain a hidden state that is updated at each time step based on the current observation and the previous hidden state, enabling the policy to accumulate information over time.

RPL~\cite{zhang2024rpl} employed a recurrent policy architecture for perceptive humanoid locomotion, where the hidden state of the RNN encodes a compressed representation of the terrain history. This hidden state enables the policy to make decisions that depend not only on the current terrain observation but also on the sequence of observations received over the past several steps. The RNN-based memory was shown to improve locomotion performance on terrains with occluded or partially observed features, where the robot must infer the terrain geometry from a sequence of limited observations.

DreamWaQ~\cite{nahraendra2024dreamwaq} also utilized a recurrent architecture, but with a specific focus on implicit terrain imagination. The recurrent hidden state in DreamWaQ serves as an implicit terrain model that encodes the essential features of the terrain encountered so far. This implicit model enables the policy to ``imagine'' terrain properties even when direct observations are unavailable, such as when the robot's sensors are temporarily occluded or when the terrain is partially outside the sensor's field of view. The key insight is that the recurrent hidden state can serve as more than just a short-term memory buffer---it can encode structured, task-relevant information about the environment that supports decision-making under partial observability.

However, recurrent memory has inherent limitations. The hidden state of an RNN has a fixed dimensionality, which limits the amount of information that can be stored. Moreover, RNNs are known to struggle with long-range dependencies: while LSTMs and GRUs mitigate the vanishing gradient problem to some extent, they still have difficulty retaining information over hundreds or thousands of time steps. For parkour scenarios that span extended courses, these limitations can result in the loss of critical information about terrain features observed early in the episode.

\subsection{Sequence Modeling and Next-Token Prediction}

A recent paradigm shift in memory-augmented locomotion has been inspired by the success of large language models (LLMs) in natural language processing. The key insight is that locomotion can be framed as a sequence prediction problem: given a sequence of observations and actions, predict the next action. This framing enables the use of powerful sequence modeling architectures---such as Transformers---that have demonstrated exceptional ability to capture long-range dependencies in sequential data.

Radosavovic et al.~\cite{radosavovic2024nexttoken} proposed Humanoid Locomotion as Next Token Prediction, a groundbreaking work that applies the Transformer architecture to humanoid locomotion control. Their method treats the sequence of proprioceptive observations and actions as a ``language'' of locomotion, and trains a Transformer model to predict the next action in this sequence, analogous to how language models predict the next token in a text sequence. The Transformer's self-attention mechanism enables the policy to attend to any previous time step in the sequence, providing effectively unlimited memory horizon. This approach demonstrated remarkable generalization: a single Transformer policy trained on diverse locomotion data could walk, run, turn, and navigate various terrains without any task-specific tuning. The success of this approach suggests that the sequence modeling paradigm, with its inherent long-range memory capabilities, is a powerful framework for locomotion control.

LocoFormer~\cite{liu2024locoformer} extended the Transformer-based locomotion paradigm with a focus on long-context adaptation. Their method employs a Transformer architecture with a long context window that can process hundreds of observation-action pairs, enabling the policy to adapt its behavior based on extensive history. The key contribution is a context adaptation mechanism that allows the policy to dynamically adjust its behavior based on the accumulated context, rather than simply processing the entire history uniformly. LocoFormer demonstrated that long-context Transformers can serve as generalist locomotion controllers that adapt to diverse terrains and tasks within a single policy, achieving state-of-the-art performance on multiple locomotion benchmarks.

The sequence modeling approach offers several advantages over recurrent memory. First, Transformers can attend to arbitrary positions in the sequence, eliminating the information bottleneck of fixed-dimensional hidden states. Second, the attention mechanism provides interpretable insights into which past observations the policy considers most relevant for the current decision. Third, the parallelizable nature of Transformer training enables more efficient computation compared to the sequential processing required by RNNs.

However, the sequence modeling approach also has limitations. The computational cost of self-attention scales quadratically with the sequence length, which can be prohibitive for very long episodes. While efficient attention variants (e.g., linear attention, sliding window attention) can mitigate this cost, they may sacrifice the ability to attend to distant time steps. Furthermore, the Transformer's reliance on a fixed context window means that information beyond the window is completely lost, which can be problematic for very long parkour courses.

\subsection{World Models and Ego-Centric Memory}

Another approach to memory-augmented locomotion leverages world models---learned models of the environment dynamics that enable the policy to simulate future outcomes and plan accordingly. World models inherently require memory, as they must maintain an internal state that represents the current belief about the environment.

Liu et al.~\cite{liu2024egovision1,liu2024egovision2} proposed an Ego-Vision World Model for humanoid contact planning, which combines a world model with ego-centric visual observations. Their method trains a generative model that predicts future visual observations given the current observation and a sequence of actions, enabling the policy to ``imagine'' the consequences of its actions before executing them. The ego-centric perspective---where observations are captured from the robot's own viewpoint---is particularly natural for contact planning, as the robot must reason about where and how its body will make contact with the environment. The world model serves as a form of predictive memory: it encodes not just what the robot has observed, but what it expects to observe in the future, enabling more informed decision-making in contact-rich scenarios.

The ego-vision world model approach has several notable properties. First, by predicting future observations in the robot's own reference frame, the model naturally handles the egocentric nature of embodied perception---the robot does not need to maintain a global map of the environment, but rather a local, action-conditioned prediction of what it will see next. Second, the generative nature of the world model enables diverse predictions, capturing the stochasticity and uncertainty inherent in real-world environments. Third, the world model can be trained in a self-supervised manner from interaction data, making it scalable to large datasets.

However, world models also introduce significant challenges. The accuracy of the model's predictions degrades over longer horizons, as small errors compound through recursive prediction. This limits the effective planning horizon of world-model-based approaches. Furthermore, training a generative world model that produces high-fidelity visual predictions is computationally expensive and requires large amounts of training data. For parkour scenarios, where the visual complexity of the environment is high and the robot must reason over extended time horizons, these challenges are particularly acute.

\subsection{Memory in the Context of Parkour}

The memory requirements of parkour present unique challenges that are not fully addressed by existing memory-augmented approaches. Parkour courses are typically long---spanning tens of meters with multiple obstacles---and the robot must maintain spatial awareness of the entire course to make effective navigation decisions. This requires memory that is both long-range (spanning hundreds of time steps) and spatially structured (encoding the geometry of the environment in a way that supports path planning and skill selection).

Existing recurrent approaches~\cite{zhang2024rpl,nahraendra2024dreamwaq} provide short-to-medium range memory but struggle with the long-range dependencies needed for extended parkour courses. The sequence modeling approaches~\cite{radosavovic2024nexttoken,liu2024locoformer} offer longer memory horizons through attention mechanisms, but their computational cost grows with sequence length and they lack explicit spatial structure that would facilitate terrain understanding. The world model approach~\cite{liu2024egovision1,liu2024egovision2} provides predictive memory but is limited by prediction accuracy over long horizons.

A particularly important gap is the lack of memory architectures that combine long-range temporal memory with spatially structured terrain representations for humanoid parkour. The ideal memory system for parkour would maintain a persistent, spatially organized representation of the terrain that the robot has observed, enabling it to recall terrain features from any point in its trajectory and use this information for planning and skill selection. This combination of temporal persistence and spatial structure is not well-served by any existing approach, motivating the development of new memory architectures specifically designed for perceptive parkour.

%% ============================================================
\section{Attention and Mixture-of-Experts in Robot Learning}
\label{sec:rw_attention}
%% ============================================================

\subsection{Attention Mechanisms: From NLP to Robotics}

The attention mechanism, originally introduced in the context of neural machine translation, has become one of the most influential architectural innovations in deep learning. At its core, attention enables a model to dynamically focus on the most relevant parts of its input, assigning different weights to different elements based on their relevance to the current task. This capability is particularly valuable in settings where the input is high-dimensional and the relevant information is sparse---a characteristic that is shared by both natural language processing and robot perception.

In the context of robot locomotion, attention mechanisms serve two primary purposes. First, they enable the policy to selectively focus on the most relevant terrain features when processing high-dimensional perceptual inputs such as heightmaps or depth images. Not all parts of the terrain are equally important for locomotion decisions: the region immediately ahead of the robot's feet is far more critical than distant terrain, and obstacles that are directly in the robot's path require more attention than those on the periphery. Attention mechanisms allow the policy to learn this selectivity automatically, rather than relying on hand-crafted feature extraction.

Second, attention mechanisms enable the policy to integrate information from multiple modalities or multiple time steps in a flexible, data-dependent manner. For example, when combining proprioceptive observations (joint angles, velocities, IMU readings) with exteroceptive observations (terrain heightmaps, depth images), attention can learn to weight each modality according to its current relevance---relying more on proprioception when the terrain is flat and predictable, and more on exteroception when the terrain is complex and uncertain.

\subsection{Attention-Based Terrain Encoding for Locomotion}

He et al.~\cite{he2024attentionmap} proposed attention-based map encoding for learning generalized legged locomotion, a pioneering work that directly applies attention mechanisms to terrain representation for locomotion control. Their method replaces the conventional convolutional neural network (CNN) encoder for heightmap processing with an attention-based encoder that can selectively focus on the most terrain-relevant regions of the heightmap.

The key insight of this work is that CNN-based terrain encoders, while effective at extracting local geometric features, have limited receptive fields that may not capture the global terrain context needed for locomotion planning. For example, a CNN with a small kernel may detect a local step or gap but fail to capture the overall terrain profile---whether the robot is approaching a staircase, navigating a ridge, or traversing a slope. Attention-based encoding, by contrast, can capture both local and global terrain features by assigning attention weights across the entire heightmap, enabling the policy to make more informed locomotion decisions.

He et al.~\cite{he2024attentionmap} demonstrated that attention-based map encoding significantly improves locomotion generalization compared to CNN-based encoding, particularly on terrains with complex global structure. The attention maps learned by the policy are interpretable: they reveal which terrain regions the policy attends to when making locomotion decisions, providing insights into the policy's decision-making process. This interpretability is a valuable property for safety-critical applications, as it enables developers to verify that the policy is attending to the correct terrain features and not relying on spurious correlations.

The attention-based approach also offers advantages in terms of parameter efficiency. Because attention mechanisms can dynamically select the most relevant features, they can achieve comparable or superior performance with fewer parameters than CNN-based encoders. This is particularly important for deployment on resource-constrained robot hardware, where computational budget is limited.

\subsection{Attention in Sequence Modeling for Locomotion}

As discussed in Section~2.2, the Transformer architecture---which is built entirely on attention mechanisms---has been applied to locomotion control through the sequence modeling paradigm. Both Humanoid Locomotion as Next Token Prediction~\cite{radosavovic2024nexttoken} and LocoFormer~\cite{liu2024locoformer} rely on self-attention to capture dependencies between observations and actions at different time steps.

In the context of sequence modeling for locomotion, attention serves a dual role. At the temporal level, self-attention enables the policy to identify and focus on the most relevant past observations when predicting the next action. For example, when the robot is about to step onto a new terrain type, the policy may attend to previous observations where it encountered similar terrain, using that experience to inform its current decision. At the feature level, attention within each observation enables the policy to focus on the most relevant components of the observation vector---for example, attending more to joint velocities when the robot is losing balance, or more to terrain features when the robot is approaching an obstacle.

The combination of temporal and feature-level attention in Transformer-based locomotion policies provides a powerful mechanism for adaptive control. However, this power comes at the cost of increased computational requirements, particularly for the self-attention computation which scales quadratically with the sequence length. This trade-off between attention capacity and computational efficiency is a central design consideration for attention-based locomotion policies.

\subsection{Mixture-of-Experts Architectures}

The Mixture-of-Experts (MoE) paradigm offers a complementary approach to attention for achieving adaptive, task-dependent processing. In an MoE architecture, the model consists of multiple ``expert'' sub-networks, each specializing in a different aspect of the input or task, and a ``gating'' network that dynamically selects which experts to activate for each input. This enables the model to allocate its computational resources adaptively---using different experts for different terrain types, locomotion modes, or skill requirements.

While MoE architectures have not been extensively explored in the locomotion literature reviewed here, the concept of expert specialization is implicitly present in several works. The skill-chaining approaches of Extreme Parkour~\cite{cheng2024extremeparkour} and Robot Parkour Learning~\cite{zhuang2024robotparkour} can be viewed as a form of MoE, where each skill is an ``expert'' and the high-level planner serves as the ``gate'' that selects the appropriate expert for the current terrain. Similarly, the implicit-explicit learning of PIE~\cite{luo2024pie} employs two parallel processing streams---implicit and explicit---that can be seen as two experts with different specializations, combined through a learned fusion mechanism.

The attention mechanism itself can be interpreted through the lens of MoE: each attention head specializes in capturing a particular type of dependency or relationship in the input, and the multi-head attention architecture combines these specializations to produce a rich, multi-faceted representation. This connection between attention and MoE suggests that hybrid architectures---combining explicit expert modules with attention-based routing---could offer the best of both worlds: the interpretability and modularity of MoE with the flexibility and data-dependence of attention.

\subsection{Attention for Multi-Modal Fusion in Locomotion}

A particularly important application of attention in locomotion is multi-modal fusion---combining information from different sensor modalities (proprioception, exteroception, terrain maps) into a unified representation for decision-making. The challenge of multi-modal fusion lies in the heterogeneity of the modalities: proprioceptive signals are low-dimensional, high-frequency, and directly related to the robot's state, while exteroceptive signals are high-dimensional, lower-frequency, and provide information about the external environment.

Bank et al.~\cite{bank2024hybrid} addressed multi-modal fusion in the context of heightmap generation, combining LiDAR and depth camera data through a hybrid autoencoder. While their method does not explicitly use attention for fusion, the challenge they address---combining heterogeneous sensor modalities with different noise characteristics and spatial resolutions---is precisely the type of problem where attention-based fusion excels. An attention-based fusion mechanism could learn to weight each modality according to its reliability and relevance for the current terrain region, producing more robust heightmaps than fixed-weight fusion.

DreamWaQ~\cite{nahraendra2024dreamwaq} implicitly performs multi-modal fusion through its recurrent hidden state, which combines proprioceptive history with implicit terrain imagination. The recurrent architecture learns to balance the influence of current proprioceptive observations against the accumulated terrain model, but this balance is implicit and may not be optimal for all situations. An explicit attention-based fusion mechanism could provide more flexible and interpretable multi-modal integration, dynamically adjusting the balance between proprioception and terrain understanding based on the current locomotion context.

\subsection{Limitations and Open Questions}

Despite the promising results of attention-based approaches in locomotion, several open questions remain.

First, the computational cost of attention remains a significant concern for real-time deployment. Locomotion policies must operate at high control frequencies (typically 50--1000\,Hz), and the quadratic complexity of self-attention can be prohibitive at these frequencies. Efficient attention mechanisms---such as linear attention, flash attention, or sparse attention---are necessary to make attention-based policies practical for real-time control, but their impact on locomotion performance has not been thoroughly evaluated.

Second, safety guarantees for attention-based locomotion policies remain an open problem. While control barrier functions (CBFs) provide formal safety guarantees for continuous control systems~\cite{tan2024distributed}, integrating such safety constraints with learned attention-based policies is challenging. CBFs enforce safety by ensuring that the system state remains within a safe set, but the high-dimensional, nonlinear dynamics of humanoid locomotion make it difficult to construct CBFs that are both valid and permissive enough to allow agile parkour maneuvers. The distributed implementation of CBFs for multi-agent systems~\cite{tan2024distributed} demonstrates that safety constraints can be enforced in a scalable manner, but extending this approach to the single-agent, high-dimensional setting of humanoid parkour---where the ``agents'' are the individual joints or body segments that must coordinate to maintain whole-body safety---remains an open research question.

Third, the interaction between attention and memory in locomotion policies is not well understood. Attention mechanisms provide a natural way to access memory (by attending to stored representations), but the design of memory systems that are optimized for attention-based access---particularly for spatially structured terrain representations---remains an open problem. The attention-based map encoding of He et al.~\cite{he2024attentionmap} processes a single heightmap observation, but extending this approach to attend over a sequence of accumulated terrain observations (i.e., a memory-augmented terrain map) has not been explored.

Fourth, the combination of attention with MoE for locomotion is largely unexplored. While the individual benefits of attention (flexible, data-dependent feature selection) and MoE (modular, scalable expert specialization) are well-established, their synergy in the context of locomotion control---particularly for parkour with diverse skill requirements---has not been investigated. A unified architecture that uses attention for terrain encoding and MoE for skill-dependent processing could offer significant advantages for humanoid parkour, but designing and training such an architecture presents substantial challenges.

%% ============================================================
\section{Adversarial Motion Priors}
\label{sec:rw_amp}
%% ============================================================

\subsection{Motion Imitation and Style Transfer: Foundations}

The goal of producing natural, human-like motion in simulated and robotic characters has been a longstanding challenge in computer graphics and robotics. Early approaches relied on tracking-based methods, where a physics-based character is controlled to follow reference motion trajectories as closely as possible. While effective for reproducing specific motions, tracking-based methods are inherently limited: they can only replicate motions that are explicitly provided as reference, and they struggle to generalize to new situations that require deviations from the reference trajectory.

DeepMimic~\cite{peng2018deepmimic} represented a major advance in motion imitation by combining deep reinforcement learning with reference motion tracking. Their method trains a policy that is rewarded for both imitating a reference motion and accomplishing a task objective (e.g., reaching a target location, performing a backflip). The key innovation is the use of a reference motion reward that encourages the simulated character to match the pose and velocity of the reference motion at each time step, while the task reward provides the motivation for goal-directed behavior. DeepMimic demonstrated that this combination of imitation and task rewards can produce highly natural and task-appropriate motions for a wide range of skills, including walking, running, jumping, and acrobatics.

However, DeepMimic has an important limitation: it requires a separate reference motion for each skill, and the quality of the learned motion is directly tied to the quality and coverage of the reference data. This means that the policy cannot produce motions that are not represented in the reference dataset, limiting its ability to generalize to novel situations or blend between different motion styles. This limitation motivates the development of methods that can learn a general notion of motion style from data, rather than requiring explicit reference motions for each skill.

\subsection{Adversarial Motion Priors (AMP)}

Adversarial Motion Priors (AMP)~\cite{peng2021amp} addressed the fundamental limitation of DeepMimic by introducing an adversarial learning framework that encodes motion style as a learned prior, rather than requiring explicit reference motion tracking. The key idea is to train a discriminator network that distinguishes between ``real'' motions (drawn from a dataset of reference motions) and ``fake'' motions (produced by the policy). The policy is then trained to generate motions that fool the discriminator, effectively learning to produce motions that are stylistically consistent with the reference dataset without explicitly tracking any single reference trajectory.

The AMP framework offers several important advantages over DeepMimic. First, it does not require temporal alignment between the policy's motion and the reference data: the discriminator evaluates motion style based on local pose statistics rather than frame-by-frame correspondence, allowing the policy to produce motions that are stylistically similar to the reference data but temporally unconstrained. This is crucial for robotics applications, where the robot's dynamics may not match the timing of the reference motion. Second, AMP can learn from a diverse dataset of reference motions spanning multiple skills and styles, enabling the policy to produce a wide range of behaviors from a single training procedure. Third, the adversarial reward provides a dense, continuous signal that guides the policy toward natural-looking motion, complementing the sparse task reward that guides the policy toward goal-directed behavior.

The AMP framework has become the standard approach for incorporating motion style into physics-based character control, and its influence extends well beyond the original computer graphics setting. In the robotics context, AMP provides a principled way to ensure that learned locomotion policies produce motions that are not only functional but also natural and energy-efficient, which is critical for real-world deployment where unnatural motions can lead to excessive energy consumption, joint wear, and instability.

\subsection{Improved Discriminator Designs}

The quality of AMP-trained policies depends critically on the design of the discriminator. A poorly designed discriminator may fail to capture important aspects of motion style, leading to policies that produce motions that are technically ``realistic'' according to the discriminator but visually or physically unsatisfactory. Several works have proposed improvements to the discriminator design.

Zhang et al.~\cite{zhang2024advdiscriminators} introduced adversarial differential discriminators for physics-based motion imitation. Their key insight is that the standard AMP discriminator evaluates motion style based on individual poses, without considering the temporal derivatives (velocities, accelerations) that are crucial for motion quality. By incorporating differential features---comparing not just the current pose but also the rate of change of the pose---the differential discriminator can more accurately distinguish between natural and unnatural motions. This leads to policies that produce smoother, more physically plausible motions, particularly for dynamic skills such as running and jumping where the temporal characteristics of the motion are as important as the spatial pose.

The differential discriminator approach also addresses a common failure mode of standard AMP: mode collapse, where the policy converges to a narrow range of motions that fool the discriminator but lack diversity. By evaluating motion style across multiple temporal scales---instantaneous pose, short-term velocity, and longer-term acceleration---the differential discriminator provides a richer training signal that encourages the policy to explore a wider range of motion styles, resulting in more diverse and adaptive behaviors.

\subsection{Motion Retargeting for Humanoid Robots}

A critical challenge in applying AMP and motion imitation to humanoid robots is the mismatch between the human motion reference data and the robot's embodiment. Human motion capture data is captured from human actors with different body proportions, joint limits, and actuation capabilities than those of a humanoid robot. Directly applying human motion as a reference for robot training can lead to infeasible or unsafe motions that exceed the robot's physical capabilities.

Araujo et al.~\cite{araujo2024retargeting} addressed this challenge with General Motion Retargeting for Humanoid Motion Tracking. Their method learns a retargeting function that maps human motion capture data to the robot's motion space, accounting for differences in body proportions, joint limits, and actuation models. The retargeting is trained jointly with the motion tracking policy, enabling end-to-end optimization that produces motions that are both faithful to the human reference and physically feasible for the robot. This work demonstrated that proper retargeting is essential for successful motion imitation learning: without retargeting, the policy may attempt to track infeasible human motions, leading to training instability and poor sim-to-real transfer.

The retargeting problem is particularly acute for humanoid parkour, where the reference motions may include extreme poses and dynamic maneuvers that push the boundaries of the robot's capabilities. A retargeting function that can adapt these motions to the robot's morphology while preserving the essential dynamics and style is crucial for enabling the robot to learn parkour skills from human demonstration.

Yang et al.~\cite{yang2024omniretarget} proposed OmniRetarget, an interaction-preserving data generation framework for humanoid whole-body loco-manipulation and scene interaction. Unlike previous retargeting methods that focus solely on kinematic correspondence, OmniRetarget preserves the interaction patterns between the human and the environment---how the hands contact objects, how the feet contact the ground, and how the body navigates through spaces. This interaction-preserving retargeting is essential for loco-manipulation tasks where the robot must not only move like a human but also interact with the environment in a human-like manner. For parkour, this means preserving the contact patterns---how the hands grip a wall, how the feet land on a narrow surface---that are critical for successful maneuver execution.

\subsection{Teleoperation and Motion Tracking}

An alternative to learning from pre-recorded motion capture data is to collect motion data directly from the robot through teleoperation. Teleoperation enables a human operator to control the robot in real time, producing motion demonstrations that are inherently feasible for the robot's embodiment because they are generated through the robot's own actuators and sensors.

TWIST~\cite{ze2024twist} proposed a Teleoperated Whole-Body Imitation System that enables whole-body motion collection for humanoid robots through teleoperation. Their system allows a human operator to control the humanoid's upper body through a VR-based interface while the lower body is controlled autonomously to maintain balance and execute locomotion. The collected teleoperation data is then used to train an imitation learning policy that can reproduce the demonstrated behaviors autonomously. TWIST demonstrated that teleoperation-based data collection can produce high-quality whole-body motions that are more feasible and natural for the robot than motions retargeted from human motion capture, because the demonstrations are generated through the robot's own kinematic and dynamic structure.

Chen et al.~\cite{chen2024gmt} proposed GMT (General Motion Tracking) for humanoid whole-body control, which trains a policy to track arbitrary reference motions in a general-purpose manner. Unlike task-specific motion tracking approaches, GMT learns a universal tracking controller that can follow any reference motion within the robot's capability, without requiring per-motion tuning. The key technical contribution is a progressive training procedure that starts with easy motions (slow walking) and gradually increases the difficulty (running, jumping, acrobatics), enabling the controller to develop the full range of tracking capabilities. GMT demonstrated that a single tracking controller can reproduce diverse human motions on a humanoid robot, providing a flexible foundation for whole-body control.

The combination of teleoperation data with AMP-style adversarial training represents a promising direction for humanoid parkour. Teleoperation provides robot-feasible reference motions, while AMP provides the style prior that enables the policy to generalize beyond the specific demonstrations. However, teleoperation of dynamic parkour maneuvers is extremely challenging, as the operator must execute precise, high-speed whole-body motions through an interface with inherent latency and limited dexterity. This limits the diversity and quality of parkour demonstrations that can be collected through teleoperation.

A complementary approach to teleoperation-based data collection is imitation learning from demonstration, where the policy learns to replicate expert behavior through supervised training on demonstration data. Ross et al.~\cite{ross2024reduction} showed that naive imitation learning suffers from compounding errors: small deviations from the expert's trajectory lead the policy into states not covered by the demonstration data, where it has no guidance and errors accumulate. They proposed DAgger, a no-regret online learning algorithm that iteratively queries the expert for corrective actions in the states visited by the learned policy, thereby covering the distribution of states the policy will encounter at deployment. This reduction of imitation learning to no-regret online learning provides theoretical guarantees on the policy's performance and has been widely adopted in robot learning. For humanoid parkour, where the state space is vast and the consequences of errors are severe, the DAgger-style approach of iteratively refining the policy with expert corrections is particularly relevant.

Luo et al.~\cite{luo2024grasping} addressed the problem of grasping diverse objects with simulated humanoids, demonstrating that whole-body humanoid control can be extended beyond locomotion to include dexterous manipulation. Their method trains a humanoid to approach, reach, and grasp a wide variety of objects using a combination of reinforcement learning and motion priors. While grasping is not directly a parkour skill, the technical challenges---coordinating whole-body movement with precise hand positioning, adapting to diverse object geometries, and maintaining balance during manipulation---are closely related to the challenges of parkour maneuvers that involve hand contact, such as wall-climbing and vaulting. The grasping framework demonstrates that AMP-style motion priors can be effectively combined with task-specific rewards for whole-body humanoid control, a principle that extends naturally to parkour.

\subsection{Beyond Imitation: Guided Diffusion and Masked Motion Inpainting}

Recent work has explored generative models beyond the adversarial framework for motion synthesis, offering new ways to produce diverse, high-quality motions for physics-based characters.

BeyondMimic~\cite{liao2024beyondmimic} proposed a framework that goes beyond simple motion mimicry by using guided diffusion to generate versatile humanoid control. Their method trains a diffusion model on a dataset of reference motions, learning to generate new motions that are stylistically consistent with the training data but not identical to any single example. The diffusion model is conditioned on task-level commands (e.g., ``walk to the target,'' ``jump over the obstacle''), enabling it to generate task-appropriate motions on demand. The generated motions are then tracked by a physics-based controller, which adapts the kinematic motion to the physical constraints of the simulated character. This two-stage approach---generative motion synthesis followed by physics-based tracking---decouples the creative motion generation from the physical execution, enabling more diverse and creative behaviors than direct RL training.

MaskedMimic~\cite{tessler2024maskedmimic} introduced unified physics-based character control through masked motion inpainting. Inspired by the masked language modeling objective used in BERT-style language models, MaskedMimic trains a policy to ``inpaint'' missing portions of a motion sequence given the observed portions. This enables flexible control at test time: the user can specify any subset of the motion (e.g., the upper body pose, the foot placement, the overall trajectory) and the policy fills in the remaining degrees of freedom in a way that is consistent with the learned motion prior. MaskedMimic demonstrated that this inpainting-based approach can unify diverse control modalities---keyframe tracking, trajectory following, goal-directed navigation---within a single policy, providing unprecedented flexibility for physics-based character control.

PARC~\cite{xu2024parc} proposed physics-based augmentation with reinforcement learning for character controllers. Their method uses physics simulation to augment a dataset of reference motions with physically plausible variations, increasing the diversity and robustness of the training data. The augmented data is then used to train an RL policy that can track the reference motions while adapting to perturbations and environmental variations. PARC demonstrated that physics-based data augmentation significantly improves the robustness of motion tracking controllers, enabling them to maintain motion quality under external disturbances and terrain variations that were not present in the original reference data.

\subsection{Learning from Video: Animal and Human Motion as Reference}

An emerging direction in motion learning is the use of video data as a source of motion reference, eliminating the need for motion capture systems or teleoperation setups.

Zhao et al.~\cite{zhao2024aggressivelocomotion} proposed learning aggressive animal locomotion skills for quadrupedal robots solely from monocular videos. Their method extracts motion features from YouTube videos of animals performing dynamic maneuvers (running, jumping, turning) and uses these features as a style prior for RL training. The key challenge is that video provides only 2D visual information, not the 3D joint positions needed for motion tracking. Their approach addresses this by using a pose estimation model to extract approximate 3D joint positions from the video, then training a policy to match the extracted motion style through an AMP-like adversarial objective. This work demonstrated that video-based motion learning can produce dynamic, animal-like locomotion on quadrupedal robots, opening the door to learning from the vast amount of animal and human motion video available online.

The video-based approach has significant implications for humanoid parkour. Parkour videos---featuring human athletes performing wall runs, vaults, precision jumps, and other dynamic maneuvers---are abundantly available online and could serve as a rich source of motion reference for training humanoid parkour policies. However, extracting useful motion information from parkour videos is challenging due to the fast dynamics, frequent occlusions, and complex camera angles. Developing robust methods for video-to-motion extraction in the parkour context remains an open problem.

\subsection{AMP in the Context of Parkour}

The application of AMP to humanoid parkour introduces several unique challenges and opportunities that are not present in the standard motion imitation setting.

First, parkour motions are highly dynamic and contact-rich, involving frequent transitions between different contact modes (feet-only, hands-and-feet, full-body contact). The AMP discriminator must be able to evaluate the style of these diverse contact modes, which requires training on a reference dataset that spans the full range of parkour motions. Collecting such a dataset---whether through motion capture, teleoperation, or video extraction---is challenging due to the physical risk and technical complexity of capturing dynamic parkour maneuvers.

Second, the adversarial training objective of AMP can conflict with the safety constraints of parkour. The discriminator rewards motions that look like the reference data, but the reference data may include motions that are at the boundary of the robot's physical capabilities. A policy that is too aggressive in matching the reference style may produce motions that exceed the robot's safety limits, leading to falls or hardware damage. Balancing the style reward from the discriminator with the safety constraints from the task reward is a delicate trade-off that requires careful tuning.

Third, the generalization capability of AMP is both a strength and a weakness for parkour. On one hand, AMP's ability to produce motions that are stylistically consistent with the reference data but not identical to any specific example enables the policy to adapt its motions to the specific terrain and obstacle configuration encountered during parkour. On the other hand, this generalization can lead to motions that are stylistically plausible but physically suboptimal---for example, a jump that looks like a parkour jump but does not achieve the necessary height or distance. Ensuring that the AMP-trained policy produces motions that are both stylistically appropriate and functionally effective is an ongoing challenge.

Several of the parkour works reviewed in Section~2.1 incorporate AMP-style adversarial training. Humanoid Parkour Learning~\cite{zhuang2024humanoidparkour} and Deep Whole-body Parkour~\cite{zhuang2024deepwholebody} both use adversarial motion priors to encourage natural-looking whole-body motion during parkour skill training. Perceptive Humanoid Parkour~\cite{wu2024perceptiveparkour} combines AMP with motion matching to ensure that the policy's motions are both stylistically consistent with human parkour demonstrations and physically feasible for the robot. Embrace Collisions~\cite{zhuang2024embrace} employs AMP to encourage contact-rich motions that leverage the robot's whole body for parkour maneuvers.

These works demonstrate that AMP is a valuable component of the parkour training pipeline, but they also reveal that AMP alone is insufficient for producing high-quality parkour behaviors. The adversarial style prior must be complemented by strong task rewards, safety constraints, and perceptive understanding to achieve reliable and effective parkour performance.

%% ============================================================
\section{Summary}
\label{sec:rw_summary}
%% ============================================================

This chapter has surveyed four major research areas relevant to target-driven, memory-attention-enhanced humanoid parkour. While each area has achieved significant progress, a critical examination reveals several important gaps that the proposed framework aims to address.

\subsection{Gaps in Humanoid Locomotion and Parkour}

The humanoid parkour literature, reviewed in Section~2.1, has demonstrated increasingly dynamic and diverse whole-body skills on real humanoid hardware. Frameworks such as Humanoid Parkour Learning~\cite{zhuang2024humanoidparkour}, Deep Whole-body Parkour~\cite{zhuang2024deepwholebody}, and Perceptive Humanoid Parkour~\cite{wu2024perceptiveparkour} have shown that RL-trained policies can execute jumping, climbing, and crawling maneuvers on complex obstacle courses. However, these works share several fundamental limitations:

\begin{enumerate}[leftmargin=2em]
\item \textbf{Reactive rather than predictive control.} Existing parkour policies process the current terrain observation and produce an immediate action, without maintaining a persistent model of the environment that supports anticipatory planning. The robot cannot prepare for upcoming obstacles or adjust its approach trajectory based on terrain observed earlier in the episode. While DreamWaQ~\cite{nahraendra2024dreamwaq} and PIE~\cite{luo2024pie} introduce implicit terrain representations, these are designed for quadrupedal locomotion and lack the spatial structure and memory horizon needed for humanoid parkour.

\item \textbf{Absence of target-driven navigation.} Current parkour systems assume that the robot follows a predefined course or receives velocity commands from an external operator. No existing framework enables the robot to autonomously navigate toward a specified target while selecting appropriate parkour skills. This target-driven capability is essential for real-world deployment, where the robot must make high-level navigation decisions---choosing paths, sequencing skills, and adapting its strategy---to reach distant goals.

\item \textbf{Limited spatial memory.} Parkour courses span tens of meters with multiple obstacles, requiring the robot to remember terrain features observed earlier in its trajectory. Current policies rely on short proprioceptive-exteroceptive histories that are insufficient for maintaining the spatial awareness needed for complex parkour navigation. The recurrent hidden states used in RPL~\cite{zhang2024rpl} and DreamWaQ~\cite{nahraendra2024dreamwaq} provide limited memory capacity, while the Transformer-based approaches of Radosavovic et al.~\cite{radosavovic2024nexttoken} and LocoFormer~\cite{liu2024locoformer} lack explicit spatial structure.
\end{enumerate}

\subsection{Gaps in Memory-Augmented Policies}

The memory architectures reviewed in Section~2.2 offer complementary strengths but none fully addresses the memory requirements of humanoid parkour:

\begin{enumerate}[leftmargin=2em]
\item \textbf{Recurrent memory}~\cite{zhang2024rpl,nahraendra2024dreamwaq} provides continuous, low-latency state updates but suffers from fixed-dimensional bottlenecks and difficulty retaining long-range dependencies. The hidden state cannot efficiently encode the spatial structure of a large terrain map.

\item \textbf{Sequence modeling}~\cite{radosavovic2024nexttoken,liu2024locoformer} provides flexible attention over long histories but incurs quadratic computational cost and lacks explicit spatial organization. The flat token sequence does not naturally represent the 2D spatial structure of terrain observations.

\item \textbf{World models}~\cite{liu2024egovision1,liu2024egovision2} provide predictive memory but suffer from compounding prediction errors over long horizons and high computational cost for visual prediction. Their ego-centric formulation does not easily support the accumulation of a global terrain map.
\end{enumerate}

The key gap is the absence of a memory architecture that combines \emph{long-range temporal persistence} with \emph{spatially structured terrain representation}, enabling the robot to maintain and query a growing map of observed terrain features throughout an extended parkour episode.

\subsection{Gaps in Attention and Mixture-of-Experts}

The attention and MoE mechanisms reviewed in Section~2.3 have shown promise for locomotion but have not been fully exploited for parkour:

\begin{enumerate}[leftmargin=2em]
\item \textbf{Attention-based terrain encoding}~\cite{he2024attentionmap} has been demonstrated for single-observation heightmap processing, but has not been extended to attend over accumulated terrain memory. The natural synergy between attention mechanisms and memory-augmented terrain maps---where attention selectively queries the most relevant terrain features from a persistent spatial representation---remains unexplored.

\item \textbf{Multi-modal fusion} via attention has not been systematically studied for parkour, where the policy must integrate proprioceptive feedback, exteroceptive terrain observations, target information, and motion style cues. Existing approaches either concatenate modalities without learned fusion~\cite{zhuang2024humanoidparkour} or use implicit recurrent fusion~\cite{nahraendra2024dreamwaq} that lacks the flexibility of explicit attention-based fusion.

\item \textbf{Mixture-of-Experts} for skill-dependent processing has not been combined with attention-based terrain encoding in a unified architecture. The skill-chaining approaches of Extreme Parkour~\cite{cheng2024extremeparkour} and Robot Parkour Learning~\cite{zhuang2024robotparkour} use hard skill selection at the planner level, while a soft MoE architecture with attention-based routing could enable smoother skill blending and more efficient parameter utilization.
\end{enumerate}

\subsection{Gaps in Adversarial Motion Priors}

The AMP framework and its extensions, reviewed in Section~2.4, provide powerful tools for learning natural-looking motions, but several gaps remain in their application to parkour:

\begin{enumerate}[leftmargin=2em]
\item \textbf{Motion reference diversity.} Existing AMP-based parkour methods~\cite{zhuang2024humanoidparkour,zhuang2024deepwholebody,wu2024perceptiveparkour} train on limited motion reference datasets that may not span the full range of parkour skills. The quality and diversity of the reference data directly limits the style diversity of the learned policy. Methods for scaling AMP to large, diverse motion datasets---potentially sourced from video~\cite{zhao2024aggressivelocomotion} or generative augmentation~\cite{xu2024parc}---have not been fully explored in the parkour context.

\item \textbf{Safety-style trade-off.} The adversarial style reward can conflict with safety constraints, particularly for dynamic parkour maneuvers that operate near the robot's physical limits. The differential discriminator of Zhang et al.~\cite{zhang2024advdiscriminators} improves motion quality but does not explicitly address the safety-style trade-off. A principled method for balancing style fidelity with safety guarantees is needed.

\item \textbf{Retargeting for parkour.} While general motion retargeting~\cite{araujo2024retargeting} and interaction-preserving retargeting~\cite{yang2024omniretarget} have been proposed, retargeting methods specifically designed for the extreme poses and contact-rich motions of parkour are lacking. Parkour motions involve unconventional body configurations (e.g., wall-climbing poses, vaulting arcs) that challenge standard retargeting assumptions.
\end{enumerate}

\subsection{The Research Gap Addressed by This Thesis}

The gaps identified above collectively point to a central research challenge: \emph{how to design a humanoid parkour policy that is target-driven, maintains spatially structured long-range memory of the terrain, uses attention to selectively process terrain and motion information, and produces natural-looking whole-body motions through adversarial motion priors.}

No existing work addresses all of these requirements simultaneously. The parkour frameworks of Section~2.1 achieve dynamic whole-body skills but lack target-driven navigation and long-range memory. The memory architectures of Section~2.2 provide various forms of temporal memory but lack the spatial structure needed for terrain accumulation. The attention mechanisms of Section~2.3 enable selective processing but have not been integrated with memory for terrain-aware parkour. The AMP methods of Section~2.4 ensure motion naturalness but do not address the perceptual and navigational challenges of parkour.

This thesis proposes a unified framework---Target-Driven Memory-Attention-Enhanced Humanoid Parkour---that addresses these gaps by integrating: (1)~a spatially structured terrain memory that accumulates observations over the entire parkour episode, (2)~an attention-based terrain encoder that selectively processes the most relevant terrain features from memory, (3)~a target-driven navigation module that guides skill selection toward a specified goal, and (4)~an adversarial motion prior that ensures natural, human-like whole-body motion. The design and evaluation of this framework are presented in the subsequent chapters.